\chapter{Herhaling afgeleiden, integralen en differentiaalvergelijkingen}
\label{ch:herhaling_afgeleiden_en_integralen}


\section{Afgeleiden}

\begin{conceptbox}[Differentiatie]
Differentiatie is het proces waarbij men de afgeleide van een gegeven functie $f(x)$ berekent t.o.v. 
de onafhankelijke variabele $x$. 
In de context van signalen en systemen is de onafhankelijke variabele doorgaans de tijd $t$, genoteerd als $f(t)$.

De afgeleide $\dd{f(t)}{t}$ of $f'(t)$ is de snelheid waarmee de functie verandert op tijdstip $t$. Meetkundig is 
dit de \textbf{helling van de raaklijn}.
\end{conceptbox}

Hieronder zie je visueel dat de rechtse functie de snelheid voorstelt waarmee de linkse functie verandert.

\begin{center}
\begin{minipage}{0.45\textwidth}
\centering
\begin{tikzpicture}
\begin{axis}[
    schoolEngineeringPlot,
    title={Functie $f(t)$},
    xlabel={$t$}, ylabel={$f(t)$},
    width=\linewidth, height=5cm,
    domain=-2.5:2.5, samples=100,
    axis lines=middle,
    xtick={0}, ytick={0}
]
\addplot[blue, thick] {x^3/3 - x};
\node[anchor=north east] at (axis description cs:0.9,0.9) {$f$};
\end{axis}
\end{tikzpicture}
\end{minipage}
\hfill
\begin{minipage}{0.45\textwidth}
\centering
\begin{tikzpicture}
\begin{axis}[
    schoolEngineeringPlot,
    title={Afgeleide $f'(t)$ (snelheid)},
    xlabel={$t$}, ylabel={$f'(t)$},
    width=\linewidth, height=5cm,
    domain=-2.5:2.5, samples=100,
    axis lines=middle,
    xtick={0}, ytick={0}
]
\addplot[green!60!black, thick] {x^2 - 1};
\node[anchor=north east] at (axis description cs:0.9,0.9) {$f'$};
\end{axis}
\end{tikzpicture}
\end{minipage}
\end{center}

\begin{theorie}[Rekenregels Afgeleiden]
\textbf{Vier fundamentele afgeleiden:}
\begin{enumerate}
    \item Machtsfunctie: $(t^n)' = n t^{n-1}$
    \item Sinusfunctie: $(\sin t)' = \cos t$
    \item Cosinusfunctie: $(\cos t)' = -\sin t$
    \item Exponentiële functie: $(e^t)' = e^t$
\end{enumerate}

\textbf{Zes differentiatie-eigenschappen:}
Als $f, g$ afleidbaar zijn:
\begin{enumerate}
    \item Constante: $c' = 0$
    \item Som: $(f + g)'(t) = f'(t) + g'(t)$
    \item Product: $(f g)'(t) = f'(t)g(t) + f(t)g'(t)$
    \item Quotiënt ($g(t) \neq 0$): $\displaystyle \left(\frac{f}{g}\right)'(t) = \frac{f'(t)g(t) - f(t)g'(t)}{g^2(t)}$
    \item \textbf{Kettingregel} (samengestelde functie): $(f \circ g)'(t) = f'(g(t)) \cdot g'(t)$
    \item Inverse functie: $(f^{-1})'(t) = \frac{1}{f'(f^{-1}(t))}$
\end{enumerate}
\end{theorie}

\begin{voorbeeld}[Afgeleide berekenen]
Bereken de afgeleide van $f(t) = e^{-4t} \cos 3t$.

\textbf{Stap 1: Productregel}
\[ f'(t) = (e^{-4t})' \cos 3t + e^{-4t} (\cos 3t)' \]

\textbf{Stap 2: Kettingregel}
\begin{itemize}
    \item $(e^{-4t})' = e^{-4t} \cdot (-4t)' = -4e^{-4t}$
    \item $(\cos 3t)' = -\sin 3t \cdot (3t)' = -3\sin 3t$
\end{itemize}

\textbf{Resultaat:}
\[ \boxed{f'(t) = -4e^{-4t} \cos 3t - 3e^{-4t} \sin 3t = -e^{-4t}(4\cos 3t + 3\sin 3t)} \]
\end{voorbeeld}

\begin{voorbeeld}[Quotiëntregel]
Bereken de afgeleide van $f(t) = \frac{\sin t}{t^2}$.

Met $u(t) = \sin t$ en $v(t) = t^2$:
\[ \left(\frac{u}{v}\right)' = \frac{u'v - uv'}{v^2} \]
\begin{itemize}
    \item $u'(t) = \cos t$
    \item $v'(t) = 2t$
\end{itemize}
Invullen:
\[ \boxed{f'(t) = \frac{(\cos t)(t^2) - (\sin t)(2t)}{(t^2)^2} = \frac{t^2 \cos t - 2t \sin t}{t^4} = \frac{t \cos t - 2 \sin t}{t^3}} \]
\end{voorbeeld}

\begin{voorbeeld}[Dubbele Kettingregel]
Bereken de afgeleide van $f(t) = \sin^3(4t) = (\sin(4t))^3$.

Dit is een kettingregel in een kettingregel: $h(g(k(t)))$ met $x^3$, $\sin x$ en $4t$.
\begin{enumerate}
    \item Buitenste macht afleiden: $3(\sin(4t))^2 \cdot (\sin(4t))'$
    \item Sinus afleiden: $3\sin^2(4t) \cdot (\cos(4t)) \cdot (4t)'$
    \item Binnenste lineaire functie afleiden: $3\sin^2(4t) \cos(4t) \cdot 4$
\end{enumerate}
\textbf{Resultaat:}
\[ \boxed{f'(t) = 12 \sin^2(4t) \cos(4t)} \]
\end{voorbeeld}

\section{Integratie}

\begin{conceptbox}[Primitieve functie \& Onbepaalde integraal]
Als $F'(t) = f(t)$, dan noemen we $F(t)$ een \concept{primitieve functie} van $f(t)$.
De \concept{onbepaalde integraal} is de verzameling van alle primitieve functies:
\[ \int f(t) \dt = F(t) + c, \quad \forall c \in \R \]
\end{conceptbox}

\begin{theorie}[Rekenregels Integralen]
\textbf{Lineariteit:}
\[ \int [a f_1(t) + b f_2(t)] \dt = a \int f_1(t) \dt + b \int f_2(t) \dt \]

\textbf{Integratie door substitutie:}
Als $f(t) = g'(t) h(g(t))$, stel dan $u = g(t)$, dan is $du = g'(t) \dt$:
\[ \int f(t) \dt = \int h(u) du \]
\end{theorie}

\begin{voorbeeld}[Substitutie voorbeeld]
Bereken $\displaystyle \int \frac{e^{1/t^2}}{t^3} \dt$.

Herschrijf de integrand: $\displaystyle \frac{1}{t^3} e^{t^{-2}}$.
Stel $u = t^{-2}$, dan is $du = -2t^{-3} \dt \implies -\frac{1}{2} du = t^{-3} \dt$.

Substitutie:
\[ \boxed{\int e^u \cdot \left(-\frac{1}{2}\right) du = -\frac{1}{2} e^u + c = -\frac{1}{2} e^{1/t^2} + c} \]
\end{voorbeeld}

\begin{voorbeeld}[Partiële Integratie]
Bereken $\displaystyle \int t e^t \dt$.

We gebruiken de formule: $\displaystyle \int u dv = u\cdot v - \int v du$.
\begin{itemize}
    \item Stel $u = t \implies du = \dt$
    \item Stel $dv = e^t \dt \implies v = e^t$
\end{itemize}
Invullen:
\[ \boxed{\int t e^t \dt = t e^t - \int e^t \dt = t e^t - e^t + c = e^t(t-1) + c} \]
\end{voorbeeld}

\begin{voorbeeld}[Partieelbreuksplitsing]
Bereken $\displaystyle \int \frac{1}{t^2 - 1} \dt$.

We splitsen de breuk: $\frac{1}{t^2-1} = \frac{1}{(t-1)(t+1)} = \frac{A}{t-1} + \frac{B}{t+1}$.
\begin{itemize}
    \item Teller gelijkstellen: $1 = A(t+1) + B(t-1)$
    \item Kies $t=1 \implies 1 = 2A \implies A = 1/2$
    \item Kies $t=-1 \implies 1 = -2B \implies B = -1/2$
\end{itemize}
De integraal wordt:
\[ \boxed{\int \left( \frac{1/2}{t-1} - \frac{1/2}{t+1} \right) \dt = \frac{1}{2} \ln|t-1| - \frac{1}{2} \ln|t+1| + c = \frac{1}{2} \ln\left|\frac{t-1}{t+1}\right| + c} \]
\end{voorbeeld}

\begin{conceptbox}[Bepaalde integraal (Oppervlakte)]
De bepaalde integraal $\int_{t_1}^{t_2} f(t) \dt$ komt overeen met de oppervlakte 
onder de curve $f(t)$ tussen $t_1$ en $t_2$.
\end{conceptbox}

\begin{center}
\begin{minipage}{0.45\textwidth}
\centering
\begin{tikzpicture}
\begin{axis}[
    schoolEngineeringPlot,
    xlabel={$t$}, ylabel={$f(t)$},
    width=\linewidth, height=5cm,
    domain=0:3, samples=100,
    axis lines=middle, xtick={1, 2}, xticklabels={$t_1$, $t_2$}, ytick=\empty
]
\addplot[name path=f, blue, thick] {0.5*x^2 + 0.5};
\path[name path=axis] (axis cs:0,0) -- (axis cs:3,0);
\addplot [
    thick,
    color=blue,
    fill=schoolyellow, 
    fill opacity=0.5
]
fill between[
    of=f and axis,
    soft clip={domain=1:2},
];
\node at (axis cs:1.5,0.8) {$\int \ge 0$};
\end{axis}
\end{tikzpicture}
\end{minipage}
\hfill
\begin{minipage}{0.45\textwidth}
\centering
\begin{tikzpicture}
\begin{axis}[
    schoolEngineeringPlot,
    xlabel={$t$}, ylabel={$f(t)$},
    width=\linewidth, height=5cm,
    domain=0:3, samples=100,
    axis lines=middle, xtick={1, 2}, xticklabels={$t_1$, $t_2$}, ytick=\empty
]
\addplot[name path=g, blue, thick] {-0.5*x^2 - 0.5};
\path[name path=axis] (axis cs:0,0) -- (axis cs:3,0);
\addplot [
    thick,
    color=blue,
    fill=schoolyellow, 
    fill opacity=0.5
]
fill between[
    of=g and axis,
    soft clip={domain=1:2},
];
\node at (axis cs:1.5,-0.8) {$\int \le 0$};
\end{axis}
\end{tikzpicture}
\end{minipage}
\end{center}

\subsection{Stap-voor-stap Oefeningen}

Hieronder worden enkele basisconcepten stapsgewijs uitgelegd aan de hand van oefeningen.

\begin{oefenblok}[Oefening: Kettingregel]
\textbf{Opgave:} Bereken de afgeleide van $f(t) = \sin(\cos(2t - 5))$.

\textbf{Stap 1: Analyseer de structuur}
Dit is een samengestelde functie van de vorm $f(g(h(t)))$.
\begin{itemize}
    \item Buitenste functie: $\sin(\dots)$
    \item Middelste functie: $\cos(\dots)$
    \item Binnenste functie: $2t - 5$
\end{itemize}

\textbf{Stap 2: Pas de kettingregel toe (van buiten naar binnen)}
De afgeleide van de sinus is de cosinus. Laat de inhoud ongemoeid, maar vermenigvuldig met de afgeleide van de inhoud.
\[ f'(t) = \cos(\cos(2t - 5)) \cdot [\cos(2t - 5)]' \]

\textbf{Stap 3: Differentieer de volgende laag}
De afgeleide van de cosinus is de min-sinus.
\[ [\cos(2t - 5)]' = -\sin(2t - 5) \cdot (2t - 5)' \]

\textbf{Stap 4: Differentieer de binnenste functie}
\[ (2t - 5)' = 2 \]

\textbf{Stap 5: Combineer alles}
\[ f'(t) = \cos(\cos(2t - 5)) \cdot [-\sin(2t - 5)] \cdot 2 \]
\[ \boxed{f'(t) = -2 \sin(2t - 5) \cos(\cos(2t - 5))} \]
\end{oefenblok}

\begin{oefenblok}[Oefening: Afgeleide]
\textbf{Opgave:} Bereken de afgeleide van $f(t) = \sqrt{1+\cos^2 t}$.

\textbf{Stap 1: Analyse}
Dit is een samengestelde functie: $g(h(k(t)))$ met:
\begin{itemize}
    \item Buitenste: $\sqrt{u} = u^{1/2}$
    \item Middelste: $1+v^2$ (waar $v=\cos t$)
    \item Binnenste: $\cos t$
\end{itemize}

\textbf{Stap 2: Differentieer wortel}
\[ f'(t) = \frac{1}{2\sqrt{1+\cos^2 t}} \cdot (1+\cos^2 t)' \]

\textbf{Stap 3: Differentieer inhoud ($1+\cos^2 t$)}
De afgeleide van 1 is 0. De afgeleide van $\cos^2 t$ is (kettingregel!):
\[ (\cos^2 t)' = 2\cos t \cdot (\cos t)' = 2\cos t \cdot (-\sin t) = -2\sin t \cos t \]

\textbf{Stap 4: Combineer}
\[ \boxed{f'(t) = \frac{1}{2\sqrt{1+\cos^2 t}} \cdot (-2\sin t \cos t) = \frac{-\sin t \cos t}{\sqrt{1+\cos^2 t}}} \]
\end{oefenblok}

\begin{oefenblok}[Oefening: Oneigenlijke integraal]
\textbf{Opgave:} Bereken $\displaystyle \int_0^{+\infty} e^{-t/2} \dt$.

\textbf{Stap 1: Definitie oneigenlijke integraal}
Vervang de bovengrens $+\infty$ door een variabele $R$ en neem de limiet.
\[ \lim_{R \to +\infty} \int_0^R e^{-t/2} \dt \]

\textbf{Stap 2: Primitieve zoeken (Substitutie)}
Stel $u = -t/2$, dan $\diff u = -1/2 \dt \implies \dt = -2 \diff u$.
\[ \int e^u (-2) \diff u = -2e^u = -2e^{-t/2} \]

\textbf{Stap 3: Vul de grenzen in}
\[ \left[ -2e^{-t/2} \right]_0^R = (-2e^{-R/2}) - (-2e^0) = -2e^{-R/2} + 2 \]

\textbf{Stap 4: Bereken de limiet}
\[ \boxed{\lim_{R \to +\infty} (-2e^{-R/2} + 2) = -2(0) + 2 = 2} \]
\end{oefenblok}

\begin{oefenblok}[Oefening: Oneigenlijke Integraal (Substitutie)]
\textbf{Opgave:} Bereken $\displaystyle \int_{-\infty}^{+\infty} \frac{2t}{(t^2+1)^2} \dt$.

\textbf{Stap 1: Primitieve zoeken}
Substitutie $u = t^2+1 \implies \diff u = 2t \dt$.
\[ \int \frac{1}{u^2} \diff u = \int u^{-2} \diff u = -u^{-1} = -\frac{1}{t^2+1} \]

\textbf{Stap 2: Bereken grenzen (Limieten)}
\[ \int_{-\infty}^{+\infty} \dots = \lim_{R \to +\infty} \left[ -\frac{1}{t^2+1} \right]_{-R}^{R} \]
\[ = \lim_{R \to +\infty} \left( \left(-\frac{1}{R^2+1}\right) - \left(-\frac{1}{(-R)^2+1}\right) \right) \]

\textbf{Stap 3: Conclusie}
Beide termen gaan naar 0 als $R \to \infty$.
\[ \boxed{0 - 0 = 0} \]
(Dit is logisch want de integrand is een oneven functie over een symmetrisch interval).
\end{oefenblok}

\subsection{Oefeningen}
\begin{enumerate}
    \item Bereken de afgeleide $f'(t)$:
    \begin{enumerate}
        \item $f(t) = \tan t$
        \item $f(t) = t^2 \cos t$
        \item $f(t) = \sin t \tan t$
        \item $f(t) = t^3 \sin t \cos t$
        \item $f(t) = \tan(5 - \sin 2t)$
        \item $f(t) = (1-2t)^{-3}$
        \item $f(t) = \sin(\cos(2t-5))$
        \item $f(t) = \sqrt{1+\cos^2 t}$
    \end{enumerate}

    \item Bereken de onbepaalde integraal $\int f(t) \dt$:
    \begin{enumerate}
        \item $f(t) = (t^3+t)^5(3t^2+1)$
        \item $f(t) = \sqrt{2t+1}$
        \item $f(t) = t^2 \cos t^3$
        \item $f(t) = \frac{(1+\sqrt{t})^{1/3}}{\sqrt{t}}$
        \item $f(t) = (1-\cos \frac{t}{2})^2 \sin \frac{t}{2}$
        \item $f(t) = \frac{1}{t^2} \sin \frac{1}{t} \cos \frac{1}{t}$
    \end{enumerate}

    \item Bereken de bepaalde integralen:
    \begin{enumerate}
        \item $\int_{-1}^1 3t^2 \sqrt{t^3+1} \dt$
        \item $\int_{\pi/4}^{\pi/2} \cot t \csc^2 t \dt$
        \item $\int_0^{\pi/2} \frac{2\sin t \cos t}{(1+\sin^2 t)^3} \dt$
        \item $\int_0^{+\infty} e^{-t/2} \dt$
        \item $\int_{-\infty}^{+\infty} \frac{1}{1+t^2} \dt$
        \item $\int_{-\infty}^{+\infty} \frac{2t}{(t^2+1)^2} \dt$
        \item $\int_{-\infty}^{+\infty} 2te^{-t^2} \dt$
        \item $\int_{-\infty}^0 e^{-|t|} \dt$
    \end{enumerate}
\end{enumerate}

\section{Differentiaalvergelijkingen}

\subsection{Beknopte theorie}

Lineaire differentiaalvergelijkingen (LDV) met constante coëfficiënten 
beschrijven veel fysische systemen. We focussen op tweede-orde systemen:
\[ y''(t) + a_1 y'(t) + a_0 y(t) = x(t) \]
De algemene oplossing is de som van de homogene en de particuliere oplossing:
\[ y(t) = y_h(t) + y_p(t) \]

\begin{conceptbox}[Intuïtie: Waarom deze som?]
Dit is het principe van \textbf{superpositie}. Je kan het vergelijken met een trillend systeem (bv. een massa aan een veer):
\begin{itemize}
    \item \textbf{Homogene oplossing $y_h(t)$:} Hoe het systeem uit zichzelf beweegt (natuurlijke respons, bv. uitdempen na een duw). Het rechterlid is 0 (geen externe kracht).
    \item \textbf{Particuliere oplossing $y_p(t)$:} Hoe het systeem reageert op de externe kracht $x(t)$ (geforceerde respons).
\end{itemize}
De totale beweging is simpelweg de som van die twee effecten.
\end{conceptbox}

\subsubsection{Homogene oplossing $y_h(t)$}
Los de karakteristieke vergelijking op: $\lambda^2 + a_1\lambda + a_0 = 0$.

\begin{theorie}[Gevallen voor $y_h(t)$]
Afhankelijk van de discriminant $D = a_1^2 - 4a_0$:
\begin{enumerate}
    \item \textbf{Twee reële wortels} ($\lambda_1 \neq \lambda_2$):
    \[ y_h(t) = C_1 e^{\lambda_1 t} + C_2 e^{\lambda_2 t} \]
    
    \item \textbf{Eén reële wortel} ($\lambda_1 = \lambda_2 = \lambda$):
    \[ y_h(t) = C_1 e^{\lambda t} + C_2 t e^{\lambda t} \]
    
    \item \textbf{Complexe wortels} ($\lambda_{1,2} = \alpha \pm j\beta$):
    \[ y_h(t) = e^{\alpha t} (C_1 \cos \beta t + C_2 \sin \beta t) \]
\end{enumerate}
\end{theorie}

\subsubsection{Particuliere oplossing $y_p(t)$}
We gebruiken de methode van \concept{onbepaalde coëfficiënten}. De vorm van $y_p(t)$ wordt gekozen op basis van het rechterlid $x(t)$.

\begin{center}
\begin{engtable}{colspec={l l}}
    Rechterlid $x(t)$ & Voorstel $y_p(t)$ \\
    $C$ (constante) & $A$ \\
    $Ct^n$ (veelterm) & $A_n t^n + \dots + A_1 t + A_0$ \\
    $Ce^{kt}$ & $A e^{kt}$ \\
    $C \cos(\omega t)$ of $C \sin(\omega t)$ & $A \cos(\omega t) + B \sin(\omega t)$ \\
\end{engtable}
\end{center}
\begin{waarschuwing}[Resonantie]
Als een term in je voorstel voor $y_p(t)$ al voorkomt in de homogene oplossing $y_h(t)$, vermenigvuldig je voorstel dan met $t$ (of $t^2$ indien nodig).
\end{waarschuwing}

\subsection{Stap-voor-stap Oefeningen}

\begin{oefenblok}[Oefening: LDV oplossen]
\textbf{Opgave:} Los op: $y'' + 7y' + 10y = 4\sin 3t$

\textbf{Stap 1: Homogene oplossing}
Karakteristieke vgl: $\lambda^2 + 7\lambda + 10 = 0$.
Discriminant $D = 49 - 40 = 9$. Wortels: $\lambda_{1,2} = \frac{-7 \pm 3}{2} \Rightarrow -2, -5$.
\[ y_h(t) = C_1 e^{-2t} + C_2 e^{-5t} \]

\textbf{Stap 2: Particuliere oplossing}
Rechterlid is goniometrisch ($\sin 3t$). $\omega=3$ komt niet voor in $\lambda$, dus geen resonantie.
Voorstel: $y_p(t) = A \cos 3t + B \sin 3t$.
Afgeleiden:
\begin{align*}
y_p' &= -3A \sin 3t + 3B \cos 3t \\
y_p'' &= -9A \cos 3t - 9B \sin 3t
\end{align*}
Invullen in DV:
\[ (-9A \cos 3t - 9B \sin 3t) + 7(-3A \sin 3t + 3B \cos 3t) + 10(A \cos 3t + B \sin 3t) = 4 \sin 3t \]
Groeperen van sinus en cosinus:
\[ \cos 3t (-9A + 21B + 10A) + \sin 3t (-9B - 21A + 10B) = 4 \sin 3t \]
\[ \cos 3t (A + 21B) + \sin 3t (B - 21A) = 4 \sin 3t \]

Stelsel oplossen:
\[
\begin{cases}
A + 21B = 0 \implies A = -21B \\
-21A + B = 4 \implies -21(-21B) + B = 4 \implies 442B = 4 \implies B = \frac{2}{221}
\end{cases}
\]
Dus $A = -\frac{42}{221}$.
\[ \boxed{y_p(t) = \frac{-42}{221} \cos 3t + \frac{2}{221} \sin 3t} \]

\textbf{Stap 3: Algemene oplossing}
\[ \boxed{y(t) = C_1 e^{-2t} + C_2 e^{-5t} - \frac{42}{221} \cos 3t + \frac{2}{221} \sin 3t} \]
\end{oefenblok}

\begin{oefenblok}[Oefening: Resonantie]
\textbf{Opgave:} Los op: $y'' + 16y = 8\sin 4t$, met $y(0)=4, y'(0)=0$.

\textbf{Stap 1: Homogene oplossing}
$\lambda^2 + 16 = 0 \implies \lambda = \pm 4j$.
\[ y_h(t) = C_1 \cos 4t + C_2 \sin 4t \]

\textbf{Stap 2: Particuliere oplossing (Resonantie!)}
Rechterlid is $8\sin 4t$ ($\omega=4$). Dit komt al voor in $y_h(t)$.
Dus voorstel vermenigvuldigen met $t$:
\[ y_p(t) = t(A \cos 4t + B \sin 4t) = At \cos 4t + Bt \sin 4t \]

Afgeleiden bepalen (productregel!):
$y_p' = A \cos 4t - 4At \sin 4t + B \sin 4t + 4Bt \cos 4t$
$y_p'' = -4A \sin 4t - 4A \sin 4t - 16At \cos 4t + 4B \cos 4t + 4B \cos 4t - 16Bt \sin 4t$
$y_p'' = -8A \sin 4t - 16At \cos 4t + 8B \cos 4t - 16Bt \sin 4t$

Invullen in DV ($y'' + 16y$):
De termen met $t$ vallen weg tegen $16y_p$ (check dit!).
Overblijvende termen:
\[ -8A \sin 4t + 8B \cos 4t = 8 \sin 4t \]
Coëfficiënten vergelijken:
$-8A = 8 \implies A = -1$
$8B = 0 \implies B = 0$
Dus: $y_p(t) = -t \cos 4t$

\textbf{Stap 3: Algemene oplossing en Beginvoorwaarden}
$y(t) = C_1 \cos 4t + C_2 \sin 4t - t \cos 4t$
$y(0) = C_1 - 0 = 4 \implies C_1 = 4$
$y'(t) = -4C_1 \sin 4t + 4C_2 \cos 4t - (\cos 4t - 4t \sin 4t)$
$y'(0) = 4C_2 - 1 = 0 \implies C_2 = 1/4$

\textbf{Eindoplossing:}
\[ \boxed{y(t) = 4 \cos 4t + \frac{1}{4} \sin 4t - t \cos 4t} \]
\end{oefenblok}

\subsection{Oefeningen}
Los de volgende differentiaalvergelijkingen op:
\begin{enumerate}
    \item $y'' + 6y' + 5y = 0$, $y(0)=4, y'(0)=0$
    \item $y'' + 6y' + 25y = 0$, $y(0)=4, y'(0)=0$
    \item $y'' + 6y' + 9y = 0$, $y(0)=2, y'(0)=0$
    \item $y' + 3y = 4e^{-t}$, $y(0)=1$
    \item $y'' + 16y = 8\sin 2t$, $y(0)=4, y'(0)=0$
    \item $y'' + 16y = 8\sin 4t$, $y(0)=4, y'(0)=0$ (Let op: resonantie!)
    \item $y'' + 6y' + 25y = 689\sin(20t)$, $y(0)=2, y'(0)=1$
    \item $y'' + 6y' + 9y = 30\sin(3t)$, $y(0)=1, y'(0)=0$
    \item $y''' + 6y' = t$, nul-beginvoorwaarden.
\end{enumerate}